% Nouveau EGA V (cote 25): archive cote 25, site page 122 (PDF page 123).
% typescript. Draft compared: none.
% AI-transcribed by Claude Opus 5.5 (claude-opus-5-5, adaptive thinking, effort high), Message Batches API, batch ega25b (msgbatch_01A7n5YuXubUPNK1zsepi8RG), 2026-11-13; not yet checked by eye.
\scanpage{25/122}
\mspage{25}

systêmes linéaires de diviseurs sans composante fixe  sur X paramétrés
par $\overset{\hand{\vee}}{P}$ , et l'ensemble Ø  des applications rationelles de X dans P, \hand{qui}
\inserted{sur $\bar{k}$ ne se factorisent par aucun hyperplan de $P_{\unsure{\bar{k}}}$.}
\inserted[typed]{lorsque les anneaux locaux de X sont factoriels.}

En effet, \deleted{\xout{X étant régulier, X en est de même de \unsure{X $X_S$} qui est lisse}
\xout{sur X, donc d'après un résultat de Auslander-Buchsbaum (qui sera repris}
\xout{au Chap III) \illegible{1} de les anneaux locaux de X\unsure{x}P sont factoriels,}
et comme codim(X-U,X) \quad 2} , il s'ensuit que tout Module inversible sur
U peut se prolonger en un Module inversible sur X, donc la condition
envisagée dans b) est automatiquement satisfaite. D'autre part, d'après
\deleted{\unsure{un Mémoire inédit lorsque n° ultérieur}}
Auslander-Buchsbaum (\deleted{qui figurera un jour au Chap III}\inserted{cf §20}) un annea\hand{/}u local
régulier est factoriel.

\deleted{La proposition 2  b)  nous montre que myennant une hypothèse}
\deleted{\inserted[typed]{(savoir D \quad Ø)}}
\deleted{supplémentaire sur un\xout{e famille} système linéaire de diviseurs D, on peut}
\deleted{lui associer canoniquement un Module inversible $\underline{L}$ sur X, prolongement}
\deleted{\inserted[typed]{où U=X-Z.}}
\deleted{canonique de $\underline{L}^{x}_{U} = f\ (\underline{O}_P(1))$ ,  On peut en donner une autre interprétation,}
\deleted{indépendante de l'hypothèse  D \quad Ø , de la façon suivante; par constructio}
\deleted{\inserted[typed]{can isomorphe à} \hand{)}}
\deleted{\xout{l'image inverse de $\underline{L}$x} ou presque $\underline{L}^{x}_{U}\otimes_S\underline{O}_P(1)$ \xout{n'}est \xout{autre que} la restric-%}
\deleted{tion $\underline{N}'$ à $U\mathrm{x}_S\overset{\hand{\vee}}{P}$  du  Module inversible $\underline{N}$  sur $X\mathrm{x}_S\overset{\hand{\vee}}{P}$  défini par D. En vertu}
\deleted{du \xout{\unsure{n°1}} n°1, cela détermine \xout{uniquement $\underline{N}$, et} $\underline{L}_U$  à isom unique près;}
\deleted{et  (comme U est dense dans X)  il existe  un Module inversible $\underline{L}$ sur X,}
\deleted{déterminé à isom unique près (do\unsure{ç}c \unsure{prolongeant|prolongeaant} $\underline{L}_U$ ) tel que  l'on ait}
\deleted{un isomorphisme  $\underline{L}\otimes_{\underline{O}_S}\underline{O}_{\unsure{P}}\ (1) \quad \underline{N}$ ;  $\underline{L}$  est donné par exemple par}
\deleted{\[ \underline{L} \quad pr_1\ (\underline{N}) \quad ; \]}
\deleted{on l'appelle  le Module inversible sur X associé au systême lin de divi-%}
\deleted{suurs D.}

\mspage{122}
